\section{Conclusiones}

\begin{itemize}
    \item \textbf{Conservación simpléctica:} El integrador de Velocity Verlet demuestra ser idóneo para el seguimiento orbital a largo plazo. La acotación en la variación del momento angular ($\Delta L \sim 10^{-13}$) verifica numéricamente la naturaleza simpléctica del algoritmo, preservando la métrica del espacio de fases y evitando derivas seculares espurias en la energía.
    \item \textbf{Eficacia de la extrapolación paramétrica:} Ante la presencia de fluctuaciones numéricas para masas bajas, el uso de una extrapolación lineal desde el régimen $[5 M_J, 11.5 M_J]$ se revela como una estrategia computacionalmente eficiente. Este método permite aislar la tendencia de la precesión de manera precisa, filtrando el ruido de corto. Con esto, se obtiene una precesión para la masa real de Júpiter de $\dot{\omega} = -1.57 \times 10^{-4} \, \text{rad/año}$.
    \item \textbf{Inestabilidad de puntos colineales:} La caracterización de $L_1$, $L_2$ y $L_3$ como puntos de silla en la superficie del pseudopotencial $U_{\text{eff}}(x,y)$ coincide con la rápida divergencia orbital bajo variaciones $\delta_i$, confirmando su inestabilidad dinámica incondicional en el marco del Problema Circular Restringido de los Tres Cuerpos (PCR3C).
    \item \textbf{Estabilización dinámica por Coriolis:} El confinamiento de las trayectorias perturbadas en $L_4$ evidencia que, pese a ser máximos del potencial efectivo, las fuerzas de inercia dependientes de la velocidad estabilizan el sistema.
\end{itemize}


